% Vietnamese translation for OI Public Library
% Source-File: IOI中国国家候选队论文/2003/林希德/林希德.doc
\documentclass[11pt]{article}
\usepackage{oipl-vn}

\title{Tìm chuỗi con lặp lớn nhất}
\author{Lin Xide}
\date{2003}

\begin{document}
\maketitle

\origtitle{Finding maximal repetition}
\sourcepath{IOI China candidate team papers / 2003 / Lin Xide / Lin Xide.doc}

\begin{abstract}
So với nhiều chủ đề thuật toán khác, xử lý chuỗi ít được trình bày trong các
tài liệu tiếng Trung cổ điển về thi tin học. Bài viết này mô tả bài toán tìm
chuỗi con lặp lớn nhất, rồi lần lượt nhắc lại hai công cụ thường dùng trong
xử lý chuỗi: cây hậu tố và KMP. Trên cơ sở đó, bài viết trình bày một thuật
toán tuyến tính để tìm chuỗi con lặp có chu kỳ lớn nhất.
\end{abstract}

\noindent\textbf{Từ khóa:} cây hậu tố, mở rộng KMP, chuỗi con tối ưu.

\section{Bài toán và nhận xét ban đầu}

\subsection{Định nghĩa}

Với chuỗi \(S\), nếu tồn tại số nguyên dương \(p\) sao cho
\[
  1\le x\le |S|-p \quad\Longrightarrow\quad S_x=S_{x+p},
\]
thì gọi \(p\) là một \term{chu kỳ} của \(S\). Khi đó
\[
  e=\frac{|S|}{p}
\]
được gọi là \term{số mũ} của \(S\), và mọi chuỗi con độ dài \(p\) của \(S\)
đều có thể xem là một \term{chuỗi sinh} của \(S\). Đặc biệt, nếu \(e\ge 2\),
ta gọi \(S\) là một chuỗi lặp kích thước \(p\).

Bài toán: cho một chuỗi \(W\) gồm các chữ cái in hoa, độ dài
\[
  1\le n\le 10^5,
\]
hãy tìm trong mọi chuỗi con của \(W\) một chuỗi lặp có chu kỳ lớn nhất. Bài
toán này được gọi là tìm \term{chuỗi con lặp lớn nhất}.

Ký hiệu \(W(u,v)\) là chuỗi con của \(W\) bắt đầu tại vị trí \(u\) và kết
thúc tại vị trí \(v\).

\subsection{Cách trực tiếp}

Thuật toán trực tiếp là duyệt chu kỳ \(p\) từ lớn xuống nhỏ, rồi với mỗi
\(p\) quét toàn bộ chuỗi để xem có hai khối liên tiếp độ dài \(p\) giống nhau
hay không. Khi tìm thấy, có thể trả lời ngay vì ta đang xét chu kỳ theo thứ tự
giảm dần.

Ý tưởng này dễ cài đặt nhưng có độ phức tạp \(O(n^2)\). Để giảm xuống tuyến
tính, ta cần hai công cụ: cây hậu tố để giới hạn dạng xuất hiện của nghiệm,
và một biến thể KMP để tính nhanh các độ dài khớp tiền tố/hậu tố.

\section{Cây hậu tố}

\subsection{Định nghĩa}

Trước hết thêm ký tự kết thúc đặc biệt vào cuối chuỗi:
\[
  W\leftarrow W\$,
\]
trong đó \(\$\) không xuất hiện ở vị trí nào trước đó. Ký tự này bảo đảm
không có hậu tố nào là tiền tố của hậu tố khác.

Cây hậu tố là một trie nén chứa tất cả các hậu tố của \(W\). Mỗi cạnh
\(E(u,v)\) được gắn bởi một cặp chỉ số
\[
  E(u,v)=(\ell,r),
\]
biểu diễn chuỗi con \(W(\ell,r)\). Với bảng chữ cái \(A,\ldots,Z\), mỗi nút
có tối đa 26 con tương ứng với ký tự đầu của nhãn cạnh.

Nếu nối các nhãn cạnh trên đường từ gốc tới nút \(x\), ta được một chuỗi
\(S_x\). Ta gọi \(S_x\) là chuỗi tương ứng của \(x\), và \(x\) là nút tương
ứng của \(S_x\). Mỗi hậu tố \(W(i,n)\) tương ứng với một lá
\(\mathit{Leaf}_i\).

\subsection{Dựng cây}

Khi chèn hậu tố \(W(i,n)\), đặt \(\mathit{Head}_i\) là tiền tố chung dài nhất
giữa \(W(i,n)\) và một hậu tố \(W(j,n)\) đã chèn trước đó, với \(j<i\). Đặt
\(\mathit{Tail}_i\) sao cho
\[
  \mathit{Head}_i+\mathit{Tail}_i=W(i,n).
\]
Nói cách khác, \(\mathit{Head}_i\) là tiền tố dài nhất của \(W(i,n)\) đã xuất
hiện trong phạm vi \(i-1\). Sau khi tìm được nút tương ứng với
\(\mathit{Head}_i\), ta chỉ cần thêm một lá mới cho phần còn lại.

Cách tìm thô là từ gốc so sánh từng ký tự của \(W(i,n)\). Nếu đi hết một
cạnh thì tiếp tục xuống nút con; nếu dừng giữa cạnh thì tách cạnh và tạo nút
mới. Cách này đúng nhưng có thể tốn \(O(n^2)\).

\subsection{Liên kết hậu tố}

McCreight dùng \term{liên kết hậu tố} để tăng tốc. Giả sử
\[
  \mathit{Head}_{i-1}=az,
\]
trong đó \(a\) là ký tự đầu. Vì \(z\) đã xuất hiện ít nhất hai lần trong phạm
vi đã xét, \(\mathit{Head}_i\) có tiền tố \(z\). Liên kết hậu tố của nút ứng
với \(az\) trỏ tới nút ứng với \(z\). Nếu \(z\) rỗng, nó trỏ về gốc.

Khi cần tìm nút của \(\mathit{Head}_i\), thay vì bắt đầu từ gốc, ta đi theo
liên kết hậu tố của nút trước, rồi chỉ tìm tiếp phần còn thiếu. Khi tạo một
nút mới, liên kết hậu tố của nó cũng được tạo ngay bằng cách đi từ liên kết
hậu tố của cha và đi xuống theo nhãn cạnh tương ứng.

Điểm then chốt trong phân tích là đại lượng
\[
  i+|\mathit{Head}_i|
\]
không giảm khi \(i\) tăng. Vì vậy toàn bộ số ký tự mà thủ tục tìm kiếm duyệt
qua là \(O(n)\). Cây hậu tố có thể được dựng trong \(O(n)\) thời gian và
\(O(n)\) bộ nhớ.

\section{KMP và mở rộng KMP}

\subsection{KMP cổ điển}

Cho chuỗi chính \(S\) và mẫu \(T\), với \(n=|S|\), \(m=|T|\). Thuật toán khớp
mẫu trực tiếp có thể tốn \(O(nm)\), vì khi gặp một ký tự khác nhau nó quay lại
quá nhiều vị trí đã biết thông tin.

KMP cải thiện điểm này bằng hàm tiền tố. Đặt \(\pi[j]\) là độ dài lớn nhất
của một tiền tố của \(T\) đồng thời là hậu tố của \(T(1,j)\), nhưng ngắn hơn
toàn bộ \(T(1,j)\). Khi đang so sánh tới \(T_j\) và \(S_i\) mà không khớp,
ta dịch mẫu sao cho phần đã khớp dài nhất còn giữ được, tức chuyển về
\(\pi[j-1]+1\) thay vì bắt đầu lại từ đầu.

Hàm \(\pi\) được tính trong \(O(m)\), và thuật toán khớp mẫu chạy trong
\(O(n+m)\).

\subsection{Mở rộng KMP}

Trong bài toán chính, ta không chỉ cần biết mẫu có xuất hiện hay không, mà cần
nhiều giá trị độ dài khớp. Định nghĩa
\[
  B_i=\operatorname{lcp}\bigl(T,S(i,n)\bigr),\qquad 1\le i\le n,
\]
tức độ dài tiền tố chung dài nhất của \(T\) và hậu tố \(S(i,n)\).

Để tính nhanh \(B_i\), trước hết tính mảng
\[
  A_i=\operatorname{lcp}\bigl(T,T(i,m)\bigr),\qquad 2\le i\le m.
\]
Đây chính là dạng thường gọi là Z-function. Giữ một vị trí \(k\) sao cho
\(k+A_k\) xa nhất trong các vị trí đã biết. Khi tính \(A_i\), đặt
\[
  \mathit{Len}=k+A_k-1,\qquad L=A_{i-k+1}.
\]
Nếu \(i\le \mathit{Len}\) và \(L<\mathit{Len}-i+1\), ta có ngay
\[
  A_i=L.
\]
Ngược lại, bắt đầu từ phần đã biết và tiếp tục so sánh từng ký tự để mở rộng
\(A_i\), rồi cập nhật \(k=i\).

Vì biên phải \(k+A_k\) chỉ tăng, tổng số lần so sánh ký tự là \(O(m)\). Mảng
\(B\) được tính tương tự bằng cách dùng mảng \(A\) làm thông tin tham chiếu;
độ phức tạp là \(O(n)\). Do đó mở rộng KMP tính toàn bộ các độ dài khớp cần
thiết trong \(O(n+m)\).

\section{Thuật toán chính}

\subsection{Chuỗi con tối ưu}

Với chuỗi con \(S=W(u,v)\) và số nguyên dương \(L\), nếu:

\begin{itemize}
  \item \(S\) là chuỗi lặp có chu kỳ \(L\);
  \item \(S\) không thể mở rộng sang trái mà vẫn có chu kỳ \(L\), tức
  \(u=1\) hoặc \(W(u-1,v)\) không còn thỏa điều kiện;
  \item \(S\) không thể mở rộng sang phải mà vẫn có chu kỳ \(L\), tức
  \(v=n\) hoặc \(W(u,v+1)\) không còn thỏa điều kiện,
\end{itemize}

thì gọi \(S\) là một \term{chuỗi con tối ưu} có chu kỳ \(L\).

Nếu tìm được tất cả chuỗi con tối ưu và chu kỳ của chúng, bài toán ban đầu
trở thành chọn chuỗi có chu kỳ lớn nhất. Ý tưởng chung là:

\begin{enumerate}
  \item tìm một chuỗi sinh \(D\) độ dài \(L\);
  \item mở rộng \(D\) sang trái và sang phải đến khi không thể mở rộng nữa;
  \item nếu chuỗi sau mở rộng có độ dài ít nhất \(2L\), ta thu được một chuỗi
  con tối ưu chu kỳ \(L\).
\end{enumerate}

\subsection{Phân rã chuỗi}

Phân rã \(W\) thành
\[
  W=U_1+U_2+\cdots+U_m
\]
theo quy tắc sau. Giả sử đã có
\[
  P=U_1+\cdots+U_{i-1},\qquad W=P+Q.
\]
Nếu ký tự đầu \(Q_1\) chưa từng xuất hiện trong phạm vi \(|P|\), đặt
\[
  U_i=Q_1.
\]
Ngược lại, đặt \(U_i\) là tiền tố dài nhất của \(Q\) đã xuất hiện trong phạm
vi \(|P|\).

Ví dụ với
\[
  W=\texttt{ABAABABAABAAB},
\]
phân rã thu được
\[
  \texttt{|A|B|A|AB|ABAAB|AAB|}.
\]

Phân rã này được tính bằng cây hậu tố: \(U_i\) hoặc là một ký tự đơn, hoặc là
\(\mathit{Head}_{|P|+1}\). Vai trò của phân rã là giới hạn mạnh vị trí của
chuỗi con tối ưu.

\subsection{Giới hạn vị trí nghiệm}

Xét một chuỗi con tối ưu \(S\), và giả sử vị trí kết thúc của nó nằm trong
mảnh \(U_i\). Có thể bỏ qua trường hợp \(S\) bắt đầu trong chính \(U_i\), vì
khi đó \(S\) đã xuất hiện trong phần \(P\) theo định nghĩa của \(U_i\), và nó
sẽ được xét ở một mảnh trước.

Gọi \(R\) là chuỗi sinh cuối cùng của \(S\), tức đoạn độ dài \(L\) ở cuối
\(S\). Bằng các lập luận phản chứng dựa trên định nghĩa phân rã, ta thu được
hai kết luận:

\begin{itemize}
  \item \(R\) không thể bắt đầu quá xa về bên phải của \(U_{i-1}\);
  \item điểm bắt đầu của \(S\) cũng không thể nằm quá xa bên trái so với
  \(U_{i-1}+U_i\).
\end{itemize}

Kết luận gọn là: nếu đặt \(U=U_i\), và đặt \(V\) là hậu tố của
\[
  P=U_1+\cdots+U_{i-1}
\]
có độ dài
\[
  |U|+2|U_{i-1}|,
\]
thì mọi chuỗi con tối ưu kết thúc trong \(U\) đều có điểm bắt đầu nằm trong
\(V\), còn điểm kết thúc nằm trong \(U\).

Đây là ý nghĩa chính của phân rã: nó biến việc tìm kiếm trên toàn bộ chuỗi
thành các bài toán cục bộ trên \(V+U\).

\subsection{Tìm chuỗi sinh}

Với cặp \(V,U\) cố định, xét chu kỳ \(L\). Vì chuỗi con tối ưu \(S\) bắt đầu
trong \(V\), kết thúc trong \(U\), và có độ dài ít nhất \(2L\), ít nhất một
trong hai điều kiện sau đúng:

\begin{itemize}
  \item phần của \(S\) nằm trong \(V\) có độ dài ít nhất \(L\);
  \item phần của \(S\) nằm trong \(U\) có độ dài ít nhất \(L\).
\end{itemize}

Hai trường hợp đối xứng nhau. Xét trường hợp thứ hai: khi đó
\[
  U(1,L)
\]
là một chuỗi sinh của \(S\). Chỉ cần mở rộng chuỗi sinh này sang trái và sang
phải để kiểm tra xem nó tạo được chuỗi lặp dài ít nhất \(2L\) hay không.

\subsection{Hàm phụ trợ}

Để mở rộng trong \(O(1)\) trung bình, định nghĩa:
\[
  \mathit{Lp}_L=\operatorname{lcp}\bigl(U,\ U(L+1,|U|)\bigr),
\]
và
\[
  \mathit{Ls}_L=\operatorname{lcs}\bigl(V,\ V+U(1,L)\bigr),
\]
trong đó \(\operatorname{lcs}\) là độ dài hậu tố chung dài nhất. Trực tiếp
tính từng giá trị có thể chậm, nhưng toàn bộ các giá trị \(\mathit{Lp}_L\) và
\(\mathit{Ls}_L\) có thể được tính bằng mở rộng KMP trên các chuỗi thích hợp.
Vì vậy chi phí bình quân cho mỗi \(L\) là \(O(1)\).

Thuật toán tổng quát:

\begin{enumerate}
  \item Dùng cây hậu tố để phân rã \(W\) thành \(U_1,\ldots,U_m\).
  \item Với mỗi \(i\ge 2\), đặt \(U=U_i\), \(V\) là hậu tố của phần trước đó
  có độ dài \(|U|+2|U_{i-1}|\).
  \item Tính các mảng mở rộng \(\mathit{Lp}\) và \(\mathit{Ls}\) cho cặp
  \(V,U\).
  \item Với mọi \(1\le L\le |U|\), dùng \(U(1,L)\) làm chuỗi sinh, mở rộng
  sang hai phía bằng \(\mathit{Lp}_L,\mathit{Ls}_L\), rồi kiểm tra độ dài sau
  mở rộng có ít nhất \(2L\) hay không.
  \item Cập nhật đáp án bằng chu kỳ lớn nhất tìm được.
\end{enumerate}

\subsection{Độ phức tạp}

Bước phân rã dùng cây hậu tố nên tốn \(O(n)\) thời gian và bộ nhớ. Với mỗi
cặp \(V,U\), các mảng phụ trợ được tính trong \(O(|V|+|U|)\). Tổng độ dài các
mảnh \(U_i\) là \(O(n)\), và do cách chọn \(V\), tổng chi phí trên mọi mảnh
cũng tuyến tính. Việc liệt kê các chu kỳ \(L\) trong mỗi \(U\) có tổng chi phí
\[
  \sum_i O(|U_i|)=O(n).
\]
Vì vậy toàn bộ thuật toán chạy trong \(O(n)\) thời gian và dùng \(O(n)\) bộ
nhớ.

\section{Kết luận}

Cây hậu tố và KMP đều là những công cụ tinh tế nhưng không xa lạ. Điều đáng
chú ý trong bài toán này là cách kết hợp chúng: cây hậu tố cho ta một phân rã
có khả năng giới hạn hình dạng của nghiệm, còn mở rộng KMP giúp kiểm tra mọi
chu kỳ cục bộ với chi phí tuyến tính. Bản thân thuật toán tìm chuỗi con lặp
lớn nhất có thể không phải kỹ thuật thường xuyên dùng nguyên dạng, nhưng cách
phối hợp nhiều công cụ xử lý chuỗi để biến một tìm kiếm bậc hai thành tuyến
tính là kinh nghiệm đáng học.

\begin{thebibliography}{9}
\bibitem{fu}
Fu Qingxiang, Wang Xiaodong.
\textit{Algorithms and Data Structures}, second edition.

\bibitem{taocp}
Donald E. Knuth.
\textit{The Art of Computer Programming}.

\bibitem{handbook}
\textit{Handbook of Computer Science and Engineering}.

\bibitem{kolpakov}
Roman Kolpakov.
\textit{Finding Maximal Repetition}.

\bibitem{dam}
\textit{Discrete Applied Mathematics}, 1989.
\end{thebibliography}

\end{document}
